\subsubsection*{Натуральные числа: аксиомы Пеано}

Вычисления с натуральными числами опираются на небольшое число исходных положений.
Такое исходное положение называют \term{аксиомой} (axiom): в выбранной теории его принимают без доказательства,
а дальнейшие утверждения из него выводят.

\begin{definition}
\term{Натуральные числа с нулём} (natural numbers including zero) образуют множество \(\mathbb N_0\).
Для его описания выделяют начальное число \(0\) и операцию \(S\),
где \(S(n)\) означает число, следующее за \(n\).
Тройка \((\mathbb N_0,0,S)\) удовлетворяет следующим аксиомам Пеано.
\end{definition}

Запись \(n\in\mathbb N_0\) означает, что \(n\) принадлежит множеству \(\mathbb N_0\).
Буквы \(m,n\) в аксиомах обозначают произвольные его элементы.

\begin{axiom}[Начальный элемент]
\(0\in\mathbb N_0\).
\end{axiom}

\begin{axiom}[Следующее число]
Для каждого \(n\in\mathbb N_0\) определено единственное следующее число \(S(n)\in\mathbb N_0\).
\end{axiom}

\begin{axiom}[У нуля нет предшественника]
Для любого \(n\in\mathbb N_0\) выполнено \(S(n)\ne0\).
\end{axiom}

\begin{axiom}[Разные числа имеют разные следующие числа]
Если \(S(m)=S(n)\), то \(m=n\).
\end{axiom}

Последняя аксиома выражает \term{инъективность} (injectivity) операции \(S\):
равенство результатов возможно только при равенстве исходных чисел.

\begin{axiom}[Принцип математической индукции]
Пусть \(P(n)\) --- утверждение, сформулированное для \(n\in\mathbb N_0\).
Если
\begin{enumerate}[label=\arabic*.,leftmargin=2em,itemsep=3pt]
\item утверждение \(P(0)\) истинно;
\item для любого \(n\in\mathbb N_0\) из истинности \(P(n)\) следует истинность \(P(S(n))\),
\end{enumerate}
то \(P(n)\) истинно для всех \(n\in\mathbb N_0\).
\end{axiom}

В доказательстве методом \term{математической индукции} (mathematical induction)
первый пункт называют \term{базой индукции} (base case),
а второй --- \term{индукционным переходом} (inductive step).
Внутри перехода истинность \(P(n)\) временно принимают как
\term{индукционное предположение} (induction hypothesis) и из неё выводят следующий случай.
Проверка нескольких примеров не заменяет переход для произвольного \(n\).

\begin{definition}
\term{Единица} (one) --- число \(1=S(0)\).
Далее обозначают \(2=S(1)\), \(3=S(2)\) и так далее.
\end{definition}

По аксиоме об отсутствии предшественника нуля \(1\ne0\).
Таким образом, привычная запись множества имеет вид
\[
\mathbb N_0=\{0,1,2,3,\ldots\}.
\]

\begin{definition}
\term{Натуральные числа} (natural numbers) в этом сборнике --- элементы множества
\[
\mathbb N=\mathbb N_0\setminus\{0\}=\{1,2,3,\ldots\}.
\]
Знак \(\setminus\) означает удаление указанных элементов: здесь из \(\mathbb N_0\) удалён ноль.
\end{definition}

Натуральные числа используют для счёта отдельных предметов.
Ноль соответствует отсутствию предметов, единица --- одному предмету.
В литературе встречается и соглашение \(0\in\mathbb N\), поэтому принятая здесь договорённость указана явно.

\subsubsection*{Как из следования получаются сложение и умножение}

\begin{definition}
\term{Рекурсивное определение} (recursive definition) задаёт начальное значение
и правило получения следующего значения из уже определённого.
Для \(a,b\in\mathbb N_0\) сложение и умножение задают правилами
\[
\begin{aligned}
a+0&=a, & a+S(b)&=S(a+b),\\
a\cdot0&=0, & a\cdot S(b)&=a\cdot b+a.
\end{aligned}
\]
\end{definition}

Например,
\[
a+1=a+S(0)=S(a+0)=S(a).
\]
Поэтому после введения сложения вместо \(S(n)\) можно писать \(n+1\).
В индукционном переходе тогда доказывают \(P(n+1)\), предполагая \(P(n)\).
Если утверждение сформулировано только для \(n\ge1\), базой служит случай \(n=1\).

Следующие два доказательства показывают, как привычные свойства действий выводятся из определений.

\begin{theorem}[Ноль слева при сложении]
Для любого \(n\in\mathbb N_0\) выполнено \(0+n=n\).
\end{theorem}

\begin{proof}
База: при \(n=0\) рекурсивное определение даёт \(0+0=0\).
Переход: предположим, что \(0+n=n\). Тогда
\[
0+S(n)=S(0+n)=S(n).
\]
Значит, утверждение верно и для следующего числа. По принципу индукции оно верно для всех \(n\in\mathbb N_0\).
\end{proof}

Теперь рекурсивное правило умножения даёт
\[
a\cdot1=a\cdot S(0)=a\cdot0+a=0+a=a.
\]

\begin{theorem}[Ассоциативность сложения]
Для любых \(a,b,c\in\mathbb N_0\)
\[
(a+b)+c=a+(b+c).
\]
\end{theorem}

\begin{proof}
Зафиксируем произвольные \(a,b\in\mathbb N_0\) и проведём индукцию по \(c\).
База \(c=0\):
\[
(a+b)+0=a+b=a+(b+0).
\]
Предположим, что \((a+b)+c=a+(b+c)\). По рекурсивному правилу сложения
\[
\begin{aligned}
(a+b)+S(c)
&=S\bigl((a+b)+c\bigr)\\
&=S\bigl(a+(b+c)\bigr)\\
&=a+S(b+c)\\
&=a+\bigl(b+S(c)\bigr).
\end{aligned}
\]
Переход доказан. Принцип индукции завершает доказательство.
\end{proof}
